% !TeX root = ../main.tex
%------------------------------------------------------------------------------
%  Appendix F -- Worked Calculation Sequences
%  Step-by-step procedures that can be followed by hand, so that software
%  output can be verified independently.
%------------------------------------------------------------------------------
\chapter{Worked Calculation Sequences}
\label{app:calculations}

Each sequence below is set out as numbered steps, with the equation, the input
required and the check to apply at that step. They are deliberately
hand-workable: the point is to be able to verify what a program returns.
Symbols follow the list of symbols; the underlying expressions are collected
in Appendix~\ref{app:closed-form} and typical values in
Appendix~\ref{app:constants}.

\section{Deriving rock mass design parameters}
\label{sec:seq-parameters}

Inputs: $\sci$, $\mi$, $\GSI$, $\Dfac$, $\Ei$, and the confining stress range
of interest. Supports Chapter~\ref{ch:rock-mass-parameters}.

\begin{enumerate}
  \item Compute the rock mass constants:
  \[
  \mb = \mi \exp\!\left(\frac{\GSI-100}{28-14\Dfac}\right),\quad
  s = \exp\!\left(\frac{\GSI-100}{9-3\Dfac}\right),
  \]
  \[
  a_{HB} = \tfrac{1}{2} + \tfrac{1}{6}\left(\exp(-\GSI/15)-\exp(-20/3)\right).
  \]
  \emph{Check:} for $\GSI = 100$ and $\Dfac = 0$ these return
  $\mb = \mi$, $s = 1$, $a_{HB} = 0.5$, recovering the intact criterion.

  \item Compute the rock mass uniaxial strength
  $\scm = \sci\,s^{a_{HB}}$, and the tensile strength
  $\sigma_{tm} = -s\sci/\mb$.
  \emph{Check:} $\scm$ should be a small fraction of $\sci$ for jointed rock
  --- typically a few per cent at $\GSI = 40$.

  \item Estimate the rock mass modulus. Where $\Ei$ is known, use the
  Hoek--Diederichs relation~\citep{hoekdiederichs2006}
  \[
  \Erm = \Ei\left[0.02 + \frac{1-\Dfac/2}
  {1+\exp\!\left((60+15\Dfac-\GSI)/11\right)}\right].
  \]
  \emph{Check:} $\Erm/\Ei$ should fall between about 0.02 and 1, rising
  with $\GSI$ and falling with $\Dfac$; at $\Dfac = 0$ the ratio is about
  0.08 at $\GSI = 30$ and about 0.73 at $\GSI = 70$.

  \item Compare against at least one independent modulus estimate --- a
  $\GSI$--$\sci$ correlation, or an $\RMR$- or $Q$-based expression --- and
  record which value is adopted and why.

  \item Fit equivalent Mohr--Coulomb parameters over the confining stress
  range that the problem actually spans, taking
  $\sigma_{3\max}$ from the tunnel depth and geometry rather than from a
  default.
  \emph{Check:} the fitted envelope should lie within a few per cent of the
  Hoek--Brown curve across that range and will diverge outside it.

  \item Record $\sci$, $\mi$, $\GSI$, $\Dfac$, the confining range and the
  date against the derivation in the design report, so that a reviewer can
  reconstruct the reasoning.
\end{enumerate}

\section{Rock--support interaction sequence}
\label{sec:seq-support-interaction}

Inputs: $\pzero$, the opening radius $a$, rock mass parameters from the
sequence above, support
properties, and the intended installation distance behind the face. Supports
Chapters~\ref{ch:ccm} and~\ref{ch:support}.

\begin{enumerate}
  \item Compute the critical support pressure $p_{cr}$. If
  $p_{cr} \le 0$, the ground is elastic and the problem is
  structural rather than one of ground pressure.

  \item Construct the ground reaction curve: compute wall convergence for a
  series of support pressures from $\pzero$ down to zero, elastic above
  $p_{cr}$ and elasto-plastic below.
  \emph{Check:} the curve must pass through $(\pzero, 0)$ and fall
  monotonically.

  \item Compute the maximum convergence $u_{\max}$ at $\pin = 0$, and the
  plastic radius $R_p$ at the anticipated equilibrium pressure.
  \emph{Check:} $R_p/a$ above about 2 indicates a squeezing problem
  (Chapter~\ref{ch:high-stress}).

  \item From the longitudinal deformation profile, compute the convergence
  $u_{s0}$ already released at the installation distance $X$.
  \emph{Check:} $u_{s0}/u_{\max}$ should be roughly 0.25 at the face and 0.7 to
  0.9 at two to three radii behind it.

  \item Construct the support characteristic curve, using the stiffness and
  capacity expressions for the actual system, offset along the convergence
  axis by $u_{s0}$.

  \item Where several elements act together, sum the stiffnesses of those
  installed at the same convergence and offset separately any installed
  later.
  \emph{Check:} do not sum capacities --- the system reaches its limit when
  the first element does.

  \item Find the intersection of the two curves. This gives the equilibrium
  pressure $p_{\text{eq}}$ and the total convergence.

  \item Compute the factor of safety as
  $FS = p_{sm}/p_{\text{eq}}$, and check the equilibrium convergence
  against the clearance envelope and any third-party movement limit.

  \item If the margin is inadequate, test in turn: stiffer support, higher
  capacity, earlier installation, or a yielding system. The construction shows
  directly which of these helps.
\end{enumerate}

\section{Underground wedge analysis}
\label{sec:seq-wedge}

Inputs: three discontinuity orientations, their shear strength parameters, the
excavation geometry and orientation, and any water pressure. Supports
Chapters~\ref{ch:structural-analysis} and~\ref{ch:ground-behaviour}; the net
constructions are in Appendix~\ref{app:stereonet}.

\begin{enumerate}
  \item Plot the three planes and the excavation surface, and apply the
  removability test.
  \emph{Check:} if the block is not removable, stop --- no stability
  calculation is required.

  \item Determine the three lines of intersection and identify the sliding
  mode: free fall, single-plane sliding, or two-plane sliding along the
  intersection.

  \item Construct the wedge geometry and compute apex height, volume and
  weight $W$.
  \emph{Check:} wedge weight scales with the cube of the span; a doubling of
  span should give roughly eight times the weight for similar geometry.

  \item Resolve $W$ into components along and normal to the sliding
  direction, and compute the normal force on each sliding plane.

  \item Add the clamping stress contribution where the block is confined,
  and subtract the water force $U$ on each plane where water pressure acts.
  \emph{Check:} for a crown wedge in a relaxed roof, take the clamping
  contribution as zero unless it can be demonstrated.

  \item Compute the resisting force as
  $\sum \left[(N-U)\tan\phi + cA\right]$ over the sliding planes, where $N$
  is the normal force on a plane, $U$ the water force on it and $A$ its area,
  and the factor of safety as resisting divided by driving force.

  \item If $FS$ is below target, compute the required bolt force
  $T$ from equilibrium at the target $FS$, and the optimum bolt
  orientation --- generally inclined to the sliding direction rather than
  normal to the surface.

  \item Repeat with the orientation data sampled from their distributions to
  obtain a distribution of $FS$ rather than a single value.
\end{enumerate}

\section{Voussoir beam check for a laminated roof}
\label{sec:seq-voussoir}

Inputs: bed thickness $t$, span $L$, unit weight, deformation modulus of the
bed, and rock compressive strength. Supports Chapter~\ref{ch:ground-behaviour}.

\begin{enumerate}
  \item Assume an initial moment arm, conventionally $z = 0.75\,t$.
  \item Compute the arch thrust from moment equilibrium of the beam under its
  self-weight, $T = wL^2/(8z)$ per unit width, where $w = \gamma t$ is the
  self-weight of the bed per unit area.
  \item Compute the elastic shortening of the arch under that thrust and the
  resulting deflection, and update the moment arm.
  \item Iterate steps 2 and 3 until the moment arm converges.
  \emph{Check:} failure to converge indicates snap-through --- the span is
  unstable.
  \item Check crushing: the compressive stress at the contact zones must not
  exceed the rock strength, allowing for the reduced contact area.
  \item Check buckling: compare the mid-span deflection at convergence with
  the bed thickness, since a solution can converge while lying close to
  snap-through and the convergence test alone gives no margin.
  \item Repeat for a range of spans to establish the maximum stable span, and
  compare against the actual span to determine the reinforcement required.
\end{enumerate}

\section{Pillar design sequence}
\label{sec:seq-pillar}

Inputs: depth, pillar and opening dimensions, and rock strength. Supports
Chapter~\ref{ch:openings}.

\begin{enumerate}
  \item Compute the extraction ratio and the tributary area pillar stress
  $\sigma_p = \gamma z\,/(1-e)$ for a regular layout, where $\gamma$ is the
  unit weight of the overburden, $z$ the depth to the pillar and $e$ the
  extraction ratio.
  \emph{Check:} tributary area overestimates stress for a small group of
  openings with stiff abutments.
  \item Where the layout is irregular or the group small, obtain pillar stress
  numerically instead.
  \item Compute the pillar width-to-height ratio $W/H$.
  \item Compute pillar strength from an empirical formula appropriate to the
  rock type, incorporating the $W/H$ effect and the size effect.
  \emph{Check:} pillar strength must be well below $\sci$; a computed strength
  approaching the intact value indicates a misapplied formula.
  \item Compute $FS$ as strength divided by stress, and compare against the
  target for the pillar's duty and design life.
  \item Check roof and floor bearing capacity against the pillar stress.
  \emph{Check:} a strong pillar on weak roof or floor may punch at a load well
  below its own strength.
  \item Assess the post-peak behaviour: for $W/H$ below about 1.5, treat
  failure as potentially violent and check the stiffness criterion of
  Chapter~\ref{ch:high-stress}.
\end{enumerate}

\section{Empirical support design sequence}
\label{sec:seq-empirical}

Inputs: classification data from Chapter~\ref{ch:classification}, excavation
dimensions and use. Supports Chapter~\ref{ch:empirical-design}.

\begin{enumerate}
  \item Compute $\RMR_{89}$ from the five basic parameters, without the
  orientation adjustment.
  \item Compute $Q = (\RQD/J_n)(J_r/J_a)(J_w/\mathrm{SRF})$, recording the
  justification for each rating --- particularly SRF, which spans several
  orders of magnitude.
  \item Compute the equivalent dimension $D_e = B/\mathrm{ESR}$, where $B$ is
  the span or the height of the excavation, whichever governs the surface
  being designed, and ESR is the excavation support ratio, taken from the
  excavation's use and required design life (Table~\ref{tab:esr}).
  \item Enter the $Q$ support chart with $D_e$ and $Q$ to obtain the support
  category, bolt length and spacing, and shotcrete thickness.
  \item Apply the orientation adjustment for the tunnel case to the rating,
  and then obtain the maximum unsupported span and stand-up time from the
  $\RMR$ chart.
  \item For walls, apply the wall factor to $Q$ before entering the chart.
  \item Compare the two recommendations.
  \emph{Check:} substantial disagreement usually indicates an inconsistent
  rating rather than a real difference, and the inputs should be revisited
  before either result is used.
  \item Confirm that the ground lies within the case base of the charts. High
  stress, squeezing, very weak rock or unusual span place it outside, and the
  analytical or numerical route governs.
\end{enumerate}

\section{Settlement and damage assessment sequence}
\label{sec:seq-settlement}

Inputs: the tunnel diameter $D$, the depth to axis $z_0$, the volume loss
$V_L$, the trough width parameter $K$, and building geometry, foundation and
condition. Supports
Chapter~\ref{ch:settlement-damage}.

\begin{enumerate}
  \item Compute the trough width $i = K z_0$, then the trough volume per
  metre of tunnel $V_s = (V_L/100)\,\pi D^2/4$ with $V_L$ in per cent, and
  only then the maximum settlement $S_{\max} = V_s / (i \sqrt{2\pi})$.
  \emph{Check:} the factor of 100 belongs in $V_s$ and nowhere else, so a
  volume loss of 0.5\,per cent enters as $V_L = 0.5$ and $V_L/100 = 0.005$.
  $S_{\max}$ emerges in metres and is converted to millimetres before the
  screening step below. For $D = 6$\,m at $z_0 = 20$\,m in stiff clay,
  $V_L = 1$\,per cent gives $S_{\max}$ of order 11\,mm.
  \item Construct the transverse Gaussian trough and, where the building is
  near the face during construction, the longitudinal profile.
  \item Compute horizontal displacement and horizontal strain, tensile at the
  trough edge and compressive over the centre.
  \emph{Check:} omitting horizontal strain will understate damage, often by a
  full category.
  \item Screen the building stock: exclude a building only where the
  predicted settlement is below about 10\,mm and the predicted slope is below
  1 in 500, both conditions being required.
  \item For the remainder, compute deflection ratio in sagging and hogging
  over the building length, treating the building as conforming to the
  greenfield trough.
  \item Compute the maximum tensile strain from bending and shear, combine it
  with the horizontal strain, and read the damage category.
  \item For buildings failing the greenfield assessment, apply relative
  stiffness modification factors and reassess.
  \emph{Check:} allow for facade openings, which reduce stiffness
  substantially below what the overall dimensions imply.
  \item For piled foundations, repeat using the subsurface movement at the
  pile toe rather than the surface settlement.
  \item Where two drives are involved, apply the modification factor to the
  second and superpose --- do not simply double the single-tunnel result.
\end{enumerate}

\section{Face support pressure sequence}
\label{sec:seq-face-pressure}

Inputs: cover, diameter, ground strength parameters, water table. Supports
Chapters~\ref{ch:ground-behaviour} and~\ref{ch:tbm}.

\begin{enumerate}
  \item Compute the water pressure at the tunnel crown, axis and invert.
  \item For cohesive ground, compute the stability ratio $N$ and compare
  against the critical value for the cover-to-diameter ratio and heading
  length.
  \item For frictional ground, compute the required effective support pressure
  from a wedge--silo analysis, using the silo reduction for the overburden
  prism.
  \emph{Check:} the required pressure should be nearly independent of cover
  once arching develops; a result rising in proportion to depth indicates the
  arching term has been omitted.
  \item Add the water pressure to obtain the minimum total support pressure,
  and add a control margin.
  \item Compute the blow-out pressure from the weight of overburden, and
  subtract a margin to obtain the maximum.
  \emph{Check:} in shallow ground the window may be narrow; if it closes,
  the alignment or the method must change.
  \item Verify that both bounds are satisfied at crown and invert
  simultaneously, allowing for the pressure gradient over the face height.
  \item Convert to machine operating parameters, stating the reference
  elevation at which the target pressure applies.
\end{enumerate}
